problem:  
Let $(M^4,g(t))$, $t\in[0,T)$, be a Ricci flow on a compact oriented four‑manifold developing a finite‑time Type‑I singularity at $T$. Suppose the singular region is connected and that parabolic blow‑ups around points of maximal curvature converge to a nonflat shrinking gradient Ricci soliton $(N^4,g_\infty,f)$.  

In four dimensions, the Chern–Gauss–Bonnet and Hirzebruch signature formulas read
\[
8\pi^2 \chi(M)=\int_M \left(\tfrac14|W|^2-\tfrac12|\mathring{\text{Ric}}|^2+\tfrac{1}{24}R^2\right)\, dV,
\qquad
12\pi^2 \tau(M)=\int_M(|W^+|^2-|W^-|^2)\, dV.
\]
  
**Problem:**  
Assuming Type‑I bounds on $|{\rm Rm}|$ and smooth convergence to $(N,g_\infty)$ on rescaled neighborhoods of the singular region, does the curvature in the singular region asymptotically accumulate the topological invariants of $M$ in the following sense?  

Does there exist a decomposition of the curvature integrals
\[
\int_M |W|^2\,dV,\qquad \int_M R^2\,dV,
\]
into contributions
\[
\int_{U_t}(\cdot)dV+\int_{M\setminus U_t}(\cdot)dV
\]
for shrinking neighborhoods $U_t$ of the singular region as $t\to T$, such that  
\[
\lim_{t\to T}\int_{U_t} \Big(\tfrac14|W|^2-\tfrac12|\mathring{\text{Ric}}|^2+\tfrac{1}{24} R^2\Big)dV
= 8\pi^2\,\chi(M)-I_\infty,
\]
where $I_\infty$ measures curvature “escaping’’ through non‑singular regions?  
In particular, can one prove that under suitable geometric isolation or integral pinching of the singularity, $I_\infty=0$, so that the singular region alone captures the entire Gauss–Bonnet–signature combination of $M$?  

Why is it a "good" problem:  
This refined question corrects prior issues by using the exact curvature integrands appearing in topological invariants—precisely those in the Chern–Gauss–Bonnet and Hirzebruch formulas—rather than ad hoc functionals. It specifically asks whether the analytic process of singularity formation can localize these curvature–topology integrals.  

It remains open and at the research frontier because:

* Analytically, it demands deep control of curvature concentration and measure convergence in Ricci flow near singularities, beyond current partial results.  
* Topologically, it tests whether global invariants like $\chi(M)$ and $\tau(M)$ can be decomposed into singular and regular contributions, offering a quantitative bridge between flow analysis and global four‑manifold topology.  
* Conceptually, if proven, it would reveal that topological information of the entire manifold becomes manifest locally at curvature blow‑up points—showing that Ricci flow “remembers’’ topology through its singularities.

Its formulation is precise, grounded in known geometric identities, and connects singular geometric analysis with deep 4‑dimensional topology, satisfying the hallmarks of a “good’’ research problem.